\chapter{Value, Quality and Low Risk}\label{ch:s1:value-quality-and-low-risk}

At the end of 2006 the US value factor had grown a dollar invested in 1963 to \$9.77. By September 2020 it had lost 57.8\% from that peak, nearly two
thirds of all its gains in forty-three years, and by July 2026 it was still 29.6\% below it. Over the same years the \omterm{def:s1:value-quality-and-low-risk:quality}{profitability factor} kept earning, 2.8\% a
year at a Sharpe ratio of 0.51, and low-beta stocks kept earning about as much as high-beta ones. Value, quality and low risk are the canon of equity
factors: the most studied, the most traded, the most argued about. This chapter builds them from point-in-time data, shows that their implementation details
decide the result, and reads their record. The build is \texttt{firm.factorlib}.

\section{Value: book-to-price and its cousins}

\begin{definition}[Book-to-price ratio, value strategy]\label{def:s1:value-quality-and-low-risk:value}
A stock's \emph{book-to-price ratio}\index{book-to-price ratio} is its book equity per share divided by its price (equivalently, book equity over market
capitalisation). A \emph{value strategy}\index{value strategy} buys stocks with high ratios of a fundamental measure (book equity, earnings, cash flow, sales)
to price and sells stocks with low ratios.
\end{definition}

Fama and French's 1992 study of the cross-section of expected returns is the standard reference for value, and their value factor, HML, is the average of the small and big value portfolios minus the average of the small and big growth portfolios, rebuilt each June with
book equity from the previous fiscal year and market equity from the previous December. Over July 1963 to July 2026 it earned 3.6\% a year at a Sharpe ratio
of 0.35 ($t = 2.77$); but 5.7\% a year up to 2006, $-5.3\%$ over 2007--2020, and 9.4\% since 2021 (\cref{fig:s1:value-quality-and-low-risk:rolling}).

\begin{omfigure}
\begin{tikzpicture}
\begin{axis}[width=11cm,height=5.2cm,xlabel={\small year},ylabel={\small 10-year mean (\% a year)},tick label style={font=\footnotesize},
  xmin=1972,xmax=2026,grid=major,grid style={gray!25},legend pos=south west,legend style={font=\scriptsize},table/col sep=comma,
  x tick label style={/pgf/number format/1000 sep={}}]
\addplot[omDef,very thick] table[x=year,y=hml_pct]{figdata/strategies-1/06-value-quality-and-low-risk/rolling.csv};
\addplot[omThm,thick] table[x=year,y=rmw_pct]{figdata/strategies-1/06-value-quality-and-low-risk/rolling.csv};
\addplot[black!50,dashed] coordinates {(1972,0) (2026,0)};
\legend{value (HML),profitability (RMW)}
\end{axis}
\end{tikzpicture}

\omcaption{The Fama--French value and \omterm{def:s1:value-quality-and-low-risk:quality}{profitability factors}: trailing 120-month mean return, annualised, at each December. Derived from the Kenneth R. French
Data Library (Fama/French 5 factors, 2x3, monthly, 202607 CRSP file); the raw series is not
redistributed.}\label{fig:s1:value-quality-and-low-risk:rolling}
\end{omfigure}

\section{Quality and profitability}

\begin{definition}[Profitability factor, quality factor]\label{def:s1:value-quality-and-low-risk:quality}
A \emph{profitability factor}\index{profitability factor} is long stocks of firms with high profitability (gross profits or operating profits over assets or
book equity) and short those with low profitability. A \emph{quality factor}\index{quality factor} generalises it to a composite of characteristics investors
should pay more for, such as profitability, growth and safety.
\end{definition}

Novy-Marx found that gross profits over assets predicts average returns about as well as book-to-market, that profitable firms earn more despite higher
valuations, and that controlling for profitability sharply improves value strategies, especially among large stocks. Fama and French's five-factor model
(2015) adds operating profitability (RMW, robust minus weak) and investment (CMA, conservative minus aggressive). Asness, Frazzini and Pedersen defined
quality as profitability, growth and safety, and found that high-quality stocks trade at higher prices but not by much, so that a quality-minus-junk factor
earns significant risk-adjusted returns in the US and 24 other countries.

The appeal of the pairing is in the correlations. From 1963 to 2026, HML and RMW correlated 0.09; over value's lost years, 2007--2020, $-0.14$. Value and
momentum correlated $-0.19$, and $-0.48$ in 2007--2020. Value and investment, on the other hand, are close relatives (0.68). Profitability kept its premium
through value's drought: 2.8\% a year and a Sharpe ratio of 0.51 over 2007--2020.

On the synthetic market, which plants no profitability premium, a profitability book (four times the latest quarterly earnings over book, as filed) loses
3.8\% a year at a Sharpe ratio of $-1.03$: it is a short-value book in disguise (its correlation with the value book is $-0.60$, and its signal's IC against
the planted value premium $-0.29$), since a high book in the denominator makes the ratio low. A characteristic measured with the same accounting number as
another factor inherits that factor's premium with the opposite sign.

\section{Low risk and betting against beta}

\begin{definition}[Low-risk anomaly, betting against beta]\label{def:s1:value-quality-and-low-risk:bab}
The \emph{low-risk anomaly}\index{low-risk anomaly} is the observation that low-beta and low-volatility stocks have earned higher risk-adjusted returns than
high-beta and high-volatility stocks: the security market line is flatter than the capital asset pricing model predicts. \emph{Betting against
beta}\index{betting against beta} exploits it with a portfolio long low-beta stocks levered to a beta of one and short high-beta stocks delevered to a beta of
one.
\end{definition}

Frazzini and Pedersen's explanation is leverage aversion: investors who cannot or will not borrow buy high-beta assets to get more market exposure, bidding
them up, so that high beta goes with low alpha. They found it in US equities, twenty other equity markets, Treasuries, corporate bonds and futures, and
found BAB's returns low when funding constraints tighten. French's portfolios sorted on beta show the flat line: the lowest beta decile, with a realised beta
of 0.59, earned 6.8\% a year over the Treasury bill; the highest, with a beta of 1.60, earned 8.6\%. The capital asset pricing model, at the market's 7.2\%
excess return, says 4.3\% and 11.5\%. A betting-against-beta portfolio built from the extreme quintiles with rolling 60-month betas earned 4.3\% a year
from 1968, Sharpe ratio 0.26 ($t = 1.95$), with a beta of $-0.06$.

\begin{omfigure}
\begin{tikzpicture}
\begin{axis}[name=a,width=6.6cm,height=5.4cm,xlabel={\small realised beta},ylabel={\small excess return (\% a year)},tick label style={font=\footnotesize},
  xmin=0.4,xmax=1.7,ymin=0,ymax=16,grid=major,grid style={gray!25},title={\small French beta deciles, 1963--2026},table/col sep=comma]
\addplot[omDef,only marks,mark=*] table[x=beta,y=ret_pct]{figdata/strategies-1/06-value-quality-and-low-risk/sml_french.csv};
\addplot[black!50,dashed,domain=0.4:1.7] {7.19*x};
\end{axis}
\begin{axis}[at={(a.east)},anchor=west,xshift=1.2cm,width=6.6cm,height=5.4cm,xlabel={\small realised beta},tick label style={font=\footnotesize},
  xmin=0.4,xmax=1.7,ymin=0,ymax=16,grid=major,grid style={gray!25},title={\small synthetic market, years 3--10},table/col sep=comma]
\addplot[omThm,only marks,mark=square*] table[x=beta,y=ret_pct]{figdata/strategies-1/06-value-quality-and-low-risk/sml_synth.csv};
\end{axis}
\end{tikzpicture}

\omcaption{Realised beta and average return of beta-sorted deciles. Left: Kenneth French's value-weighted portfolios formed on beta, excess returns over
Treasury bills, with the line the capital asset pricing model implies at the market's 7.2\% (derived statistics; the raw series is not redistributed).
Right: the synthetic market's trailing-beta deciles, total returns. Data: \texttt{s1\_fetch\_value.py}, \texttt{s1\_factors.sml}.}\label{fig:s1:value-quality-and-low-risk:sml}
\end{omfigure}

The synthetic market was built the other way: its expected returns are beta times a market premium plus the planted anomalies, and nothing rewards low beta.
Its security market line is steep (the lowest trailing-beta decile earned 2.4\% a year at a realised beta of 0.49, the highest 13.4\% at 1.39), and a
betting-against-beta book loses 5.4\% a year. A strategy whose whole premise is a flat security market line has nothing to earn where the line is not flat.

\section{Implementation details that decide the result}

\begin{center}
\small
\begin{tabular}{@{}lcccc@{}}
\toprule
synthetic value books, years 3 to 10 & Sharpe & return & beta & turnover\\
\midrule
filed book, today's price, monthly & $-0.07$ & $-0.7\%$ & $-0.11$ & 3.2\\
filed book, today's price, yearly & 0.21 & 1.9\% & $-0.12$ & 2.5\\
Fama--French timing, yearly & 0.24 & 2.2\% & $-0.12$ & 2.5\\
look-ahead book from its fiscal end, monthly & $-0.06$ & $-0.6\%$ & $-0.11$ & 3.2\\
value and momentum composite, monthly & 0.44 & 4.0\% & 0.01 & 5.1\\
\bottomrule
\end{tabular}
\end{center}

The same planted value premium gives Sharpe ratios from $-0.07$ to 0.24 depending on how book-to-price is timed. The most accurate signal (today's price,
whose rank correlation with the planted premium is 1.00) does worst, because today's price folds in the recent return: a stock that has just fallen looks
cheap, and a book rebuilt every month on today's price is short the planted momentum (its correlation with the momentum book is $-0.56$). The Fama--French
convention (a book at least six months old, divided by the price of its own date) carries less of that momentum bet and does best alone. Asness and Frazzini made the same
argument on real data from the other side: timely prices forecast true book-to-price better, and value portfolios built on them earned alphas of 3.05 to 3.78
percentage points a year against the standard method. On the synthetic market the combination is what pays: a composite of
current-price value and momentum reaches a Sharpe ratio of 0.44, better than either. The look-ahead book, using each book value before it was filed, gains
nothing here, because the synthetic book moves slowly; on real data a look-ahead of this kind is the most common way a value backtest flatters itself.

\section{Combining the canon, and factor timing}

\begin{definition}[Factor timing]\label{def:s1:value-quality-and-low-risk:timing}
\emph{Factor timing}\index{factor timing} varies the weights of a portfolio's factors over time according to forecasts of their returns, for instance from
the spread in valuations between a factor's long and short sides, its recent performance or its volatility.
\end{definition}

The case for combining the canon is diversification: value, momentum and profitability each have lost decades, and their correlations are low or negative.
The case for timing it is the temptation of the value spread: when cheap stocks are unusually cheap relative to expensive ones, value's future return
should be higher. Asness, Frazzini and Pedersen report the quality version of that: a low price of quality predicted a high future return of QMJ. The
difficulty is that timing signals are slow, their samples short, and the record of a decade like 2007--2020, when value stayed cheap and kept losing, is the
reason most practitioners hold the canon in near-fixed proportions and scale it by volatility (chapter~5) rather than time it.

\section{Strategy files}

\begin{strategyfile}{Book-to-price value}\label{strat:s1:value-quality-and-low-risk:value}
\sfield{Who pays you, and why} Investors who overreact to bad news and overpay for growth, or a risk premium for distress; the literature disagrees.
\sfield{Instruments and venues} Stocks, long cheap and short expensive, or long-only tilts.
\sfield{Signal} Book equity over market capitalisation, point in time.
\sfield{Sizing and execution} Size-neutral sorts (HML's 2x3) or characteristic weights; rebuilt yearly or monthly.
\sfield{Costs} Low turnover: two to three times the book a year here.
\sfield{How it dies} Long droughts (2007--2020); growth-led markets; accounting that no longer measures value (intangibles).
\sfield{Horizon, capacity, infrastructure} Years; very large capacity; point-in-time fundamentals.
\sfield{Backtest honestly} Book values from their filing dates; the timing of price stated; the lost decade in the sample.
\sfield{Sources} Fama and French (1992); French data: HML 3.6\% a year 1963--2026, Sharpe ratio 0.35, a 57.8\% drawdown from December 2006 to September 2020;
Asness and Frazzini (2013).
\end{strategyfile}

\begin{strategyfile}{Composite value}\label{strat:s1:value-quality-and-low-risk:composite}
\sfield{Who pays you, and why} As for value, with the measurement error of any single ratio averaged away.
\sfield{Instruments and venues} As for value.
\sfield{Signal} The average z-score of several ratios to price (book, earnings, cash flow, sales), and often momentum alongside.
\sfield{Sizing and execution} Sorts or characteristic weights on the composite.
\sfield{Costs} As for value; more with momentum in the composite.
\sfield{How it dies} As for value.
\sfield{Horizon, capacity, infrastructure} Years; several fundamentals, point in time.
\sfield{Backtest honestly} The composite's ingredients fixed in advance, not chosen on the backtest.
\sfield{Sources} This chapter's simulation (value with momentum: Sharpe ratio 0.44 against 0.24 for value alone); Asness and Frazzini (2013).
\end{strategyfile}

\begin{strategyfile}{Gross profitability}\label{strat:s1:value-quality-and-low-risk:profitability}
\sfield{Who pays you, and why} Investors who underprice durable profitability, or a premium the literature has not settled.
\sfield{Instruments and venues} Stocks, long profitable and short unprofitable firms.
\sfield{Signal} Gross profits (revenue minus cost of goods sold) over total assets.
\sfield{Sizing and execution} Sorts, often combined with value.
\sfield{Costs} Low turnover.
\sfield{How it dies} Crowding; a denominator shared with value that turns it into a disguised value bet.
\sfield{Horizon, capacity, infrastructure} Years; large capacity; income statements point in time.
\sfield{Backtest honestly} Accounting data from filing dates; its correlation with the book's other factors.
\sfield{Sources} Novy-Marx (2013); French data: RMW 3.1\% a year 1963--2026, Sharpe ratio 0.39.
\end{strategyfile}

\begin{strategyfile}{Quality minus junk}\label{strat:s1:value-quality-and-low-risk:qmj}
\sfield{Who pays you, and why} Investors who pay too little for safe, profitable, growing firms.
\sfield{Instruments and venues} Stocks in many countries.
\sfield{Signal} A composite of profitability, growth and safety scores.
\sfield{Sizing and execution} Long quality, short junk; neutral to size and industry.
\sfield{Costs} Low turnover.
\sfield{How it dies} When quality becomes expensive; the price of quality predicts its return.
\sfield{Horizon, capacity, infrastructure} Years; many accounting inputs.
\sfield{Backtest honestly} Every input point in time; the composite's recipe fixed in advance.
\sfield{Sources} Asness, Frazzini and Pedersen (2019): significant risk-adjusted returns in the US and 24 other countries.
\end{strategyfile}

\begin{strategyfile}{Betting against beta}\label{strat:s1:value-quality-and-low-risk:bab}
\sfield{Who pays you, and why} Leverage-constrained investors who buy high-beta assets instead of borrowing.
\sfield{Instruments and venues} Stocks, bonds and futures, with leverage on the long side.
\sfield{Signal} Estimated beta.
\sfield{Sizing and execution} Long low beta levered to one, short high beta delevered to one.
\sfield{Costs} Financing of the levered long side; shorting high-beta names.
\sfield{How it dies} Tightening funding (BAB's returns are low then); a steep security market line.
\sfield{Horizon, capacity, infrastructure} Months; beta estimates, financing.
\sfield{Backtest honestly} Betas known before each month; financing costs; the long side's leverage.
\sfield{Sources} Frazzini and Pedersen (2014); French beta portfolios: 4.3\% a year from 1968, Sharpe ratio 0.26, before financing and trading costs.
\end{strategyfile}

\begin{strategyfile}{Low-volatility long-only}\label{strat:s1:value-quality-and-low-risk:lowvol}
\sfield{Who pays you, and why} As for BAB, taken as a long-only portfolio with less market risk.
\sfield{Instruments and venues} Stocks, long only.
\sfield{Signal} Low beta or low volatility.
\sfield{Sizing and execution} The least risky quintile or a minimum-variance portfolio (Book~7, chapter~26).
\sfield{Costs} Low turnover.
\sfield{How it dies} Crowding into the same defensive names; rising rates hitting bond-like stocks.
\sfield{Horizon, capacity, infrastructure} Years; a risk model.
\sfield{Backtest honestly} Judge it on risk-adjusted return and against the market with its own beta.
\sfield{Sources} French beta portfolios: the lowest decile earned 6.8\% a year over bills at a beta of 0.59, 1963--2026.
\end{strategyfile}

\section{Tutorial: the lost decade of value}\label{tut:s1:value-quality-and-low-risk:tut}

\begin{tutorial}
\textbf{Goal.} Build value, profitability and betting-against-beta books from point-in-time fundamentals, time book-to-price four ways, and read the French
record. \textbf{End state:} the implementation table and \cref{fig:s1:value-quality-and-low-risk:sml}.
\begin{tutsteps}
\item \textbf{Book-to-price, timed}: with today's price or with the book's own date's price.
\omcode{code/firm/factorlib/firm_factorlib.py}{86}{96}{Book-to-price from filed book values.}\label{lst:s1:value-quality-and-low-risk:bp}
\item \textbf{\omterm{def:s1:value-quality-and-low-risk:bab}{Betting against beta}}: each side scaled to a beta of one.
\omcode{code/firm/factorlib/firm_factorlib.py}{70}{83}{The betting-against-beta book.}\label{lst:s1:value-quality-and-low-risk:bab}
\item \textbf{Run} \texttt{s1\_factors.run(name)} for the seven books, \texttt{sml()}, \texttt{correlations()}, then \texttt{s1\_fetch\_value.py} once and
\texttt{fig\_factors.py}.
\end{tutsteps}
\textbf{What to change next.} Value-weight the books (the \texttt{cap} argument of \texttt{sort\_book}); add the value spread as a timing signal and test it on the French
sample before and after 2007; plant a profitability premium and check that the profitability book finds it.
\end{tutorial}

\section{Build: the factor library}\label{bld:s1:value-quality-and-low-risk:factorlib}

\begin{build}
\bfield{Purpose} The equity factor canon from point-in-time characteristics, in the forms practitioners and papers use.
\bfield{Interface} \texttt{sort\_book(signal, universe, q, cap)}, \texttt{char\_book(signal, universe)}, \texttt{composite(*signals)},
\texttt{bab\_book(beta, universe, q)}, \texttt{book\_to\_price(anchor, cum, listed, current)}.
\bfield{Rules} Characteristics from their filing dates; the timing of prices explicit; each factor reported with its correlations to the others.
\bfield{Acceptance tests} \texttt{code/firm/factorlib/tests/}: sort and characteristic weights by hand; composites with missing values; BAB's zero beta;
book-to-price with current and fixed prices.
\bfield{Stretch} Size-neutral 2x3 sorts with NYSE breakpoints; gross profitability from income statements (Book~7, chapter~11's \texttt{firm.fundpit});
value-spread timing.
\end{build}

\begin{omsources}
\item E.~F.~Fama and K.~R.~French, ``The cross-section of expected stock returns'', \emph{Journal of Finance} 47(2), 1992.
\item R.~Novy-Marx, ``The other side of value: the gross profitability premium'', \emph{Journal of Financial Economics} 108(1), 2013.
\item C.~S.~Asness, A.~Frazzini and L.~H.~Pedersen, ``Quality minus junk'', \emph{Review of Accounting Studies} 24(1), 2019.
\item A.~Frazzini and L.~H.~Pedersen, ``Betting against beta'', \emph{Journal of Financial Economics} 111(1), 2014.
\item E.~F.~Fama and K.~R.~French, ``A five-factor asset pricing model'', \emph{Journal of Financial Economics} 116(1), 2015.
\item C.~Asness and A.~Frazzini, ``The devil in HML's details'', \emph{Journal of Portfolio Management} 39(4), 2013.
\item Kenneth R. French Data Library: Fama/French 5 factors (2x3), momentum factor, portfolios formed on beta.
\end{omsources}

\section{Exercises}

\begin{exercise}[$\star$]\label{exo:s1:value-quality-and-low-risk:1}
HML's cumulative wealth was \$9.77 per dollar at its December 2006 peak and then fell 57.8\%. What share of its cumulative gain did it lose?
\end{exercise}

\begin{exercise}[$\star$]\label{exo:s1:value-quality-and-low-risk:2}
A BAB book's low-beta side has a beta of 0.6 and its high-beta side 1.5. How many dollars does it hold long and short per dollar of beta on each side, and
what is its net dollar position?
\end{exercise}

\begin{exercise}[$\star$]\label{exo:s1:value-quality-and-low-risk:3}
The synthetic market pays 3\% a year per standard deviation of book-to-price. What does a decile book earn if its long and short deciles average $\pm 1.75$
standard deviations, with half a dollar on each side?
\end{exercise}

\begin{exercise}[$\star\star$]\label{exo:s1:value-quality-and-low-risk:4}
Explain why book-to-price with today's price is negatively correlated with momentum, and why that makes value and momentum good partners.
\end{exercise}

\begin{exercise}[$\star\star$]\label{exo:s1:value-quality-and-low-risk:5}
Why does the synthetic profitability book lose money when no profitability premium is planted?
\end{exercise}

\begin{exercise}[$\star\star$]\label{exo:s1:value-quality-and-low-risk:6}
At the market's 7.2\% excess return, what does the capital asset pricing model predict for betas of 0.594 and 1.604? Compare with the French deciles.
\end{exercise}

\begin{exercise}[$\star\star\star$]\label{exo:s1:value-quality-and-low-risk:7}
\emph{Coding.} Run \texttt{run('composite')} and \texttt{correlations()}. Explain the composite's Sharpe ratio from the two books' correlation.
\end{exercise}

\begin{exercise}[$\star\star\star$]\label{exo:s1:value-quality-and-low-risk:8}
\emph{Find the flaw.} ``Value has not worked since 2007, so it has been arbitraged away; we are dropping it.''
\end{exercise}

\section{Problem: The Lost Decade of Value}

\begin{problem}[{Weekend problem --- the canon, built and read}]\label{pb:s1:value-quality-and-low-risk:1}
The French factor data and the chapter's synthetic factor books.

\medskip\noindent\textbf{Part I --- Value.}
\begin{enumerate}
\item Define book-to-price and a \omterm{def:s1:value-quality-and-low-risk:value}{value strategy}, and describe HML's construction.
\item Give HML's return and Sharpe ratio over the whole sample and in the three periods.
\item Describe the 2007--2020 drawdown.
\item Why is the timing of the price part of the definition?
\end{enumerate}

\medskip\noindent\textbf{Part II --- Quality.}
\begin{enumerate}[resume]
\item Define the profitability and \omterm{def:s1:value-quality-and-low-risk:quality}{quality factors} and give their published results.
\item Give the correlations of HML with RMW, momentum and CMA, in the whole sample and in 2007--2020.
\item What did profitability earn during value's drought?
\item Why does the synthetic profitability book lose?
\end{enumerate}

\medskip\noindent\textbf{Part III --- Low risk.}
\begin{enumerate}[resume]
\item Define the \omterm{def:s1:value-quality-and-low-risk:bab}{low-risk anomaly} and \omterm{def:s1:value-quality-and-low-risk:bab}{betting against beta}.
\item What does French's beta-sorted data show, against the capital asset pricing model?
\item Give the BAB portfolio's return, Sharpe ratio and beta.
\item Why does BAB lose on the synthetic market?
\end{enumerate}

\medskip\noindent\textbf{Part IV --- The verdict.}
\begin{enumerate}[resume]
\item State the \emph{named result}: the value factor's drawdown from its 2007 peak in the French data, and the correlation of value and profitability returns.
\item Rank the four timings of synthetic book-to-price and explain the ranking.
\item What does combining value with momentum do here, and why?
\item What did Asness and Frazzini find about timely prices?
\item Define \omterm{def:s1:value-quality-and-low-risk:timing}{factor timing} and give one argument for and one against.
\item What would you check before trusting a value backtest?
\item How would you allocate across the canon?
\item In one sentence: what do value, quality and low risk have in common?
\end{enumerate}
\end{problem}

\section{Interview questions}

\begin{interviewq}[$\star$ \iqroles{researcher}]\label{iq:s1:value-quality-and-low-risk:1}
How is the Fama--French value factor constructed?
\end{interviewq}

\begin{interviewq}[$\star\star$ \iqroles{researcher}]\label{iq:s1:value-quality-and-low-risk:2}
Why might low-beta stocks earn higher risk-adjusted returns than high-beta stocks?
\end{interviewq}

\begin{interviewq}[$\star\star$ \iqroles{researcher, risk}]\label{iq:s1:value-quality-and-low-risk:3}
Value lost more than half its value from 2007 to 2020. What would you conclude, and what would you not?
\end{interviewq}

\begin{interviewq}[$\star\star$ \iqroles{researcher}]\label{iq:s1:value-quality-and-low-risk:4}
Why combine value with profitability, or with momentum?
\end{interviewq}

\begin{interviewq}[$\star\star$ \iqroles{developer, researcher}]\label{iq:s1:value-quality-and-low-risk:5}
What point-in-time issues arise in building a book-to-price signal?
\end{interviewq}

\begin{interviewq}[$\star\star\star$ \iqroles{researcher}]\label{iq:s1:value-quality-and-low-risk:6}
Two factors each have a Sharpe ratio of 0.3 and a correlation of $-0.5$. What is the Sharpe ratio of their equal-weighted combination, and of the best
combination?
\end{interviewq}
